\documentclass[11pt]{article}

\usepackage{amsmath, amssymb, amsthm, amsfonts, mathtools}
\usepackage[hidelinks]{hyperref}
\usepackage{xcolor}
\usepackage{xspace}
\usepackage{geometry}
\geometry{a4paper, margin=1in}

\theoremstyle{plain}
\newtheorem{theorem}{Theorem}
\newtheorem{lemma}{Lemma}

\theoremstyle{remark}
\newtheorem{remark}{Remark}

\title{Proof of Conjecture~3: The Period Length at $\omega = 1$}
\author{Dirk Kunert \thanks{Independent Researcher, Email: dirk.kunert@gmail.com}}
\date{\today}

\begin{document}

\maketitle

% =====================
\section*{Statement}
% =====================

\begin{theorem}
Let $\alpha, \beta \in \mathbb{N}$ with $\gcd(\alpha, \beta) = 1$, and let $\omega = 1$.
Then the period length of the distance sequence $\left(\delta^{(i)}\right)$ of the
cut-and-project sequence is
\[
    \lambda_{\alpha,\beta} \;=\; \Lambda_{\alpha,\beta} \;:=\; \alpha + \beta + 1.
\]
\end{theorem}

% =====================
\section*{Proof}
% =====================

\subsection*{Step 1: The window constraint as a single invariant}

For $\omega = 1$ the valid lattice points $(x, y) \in \mathbb{Z}^2$ are those satisfying
\[
    \frac{\alpha}{\beta}(x - 1) \;\le\; y \;\le\; \frac{\alpha}{\beta}\,x + 1.
\]
Multiplying through by $\beta > 0$ gives the equivalent condition
\[
    \alpha x - \alpha \;\le\; \beta y \;\le\; \alpha x + \beta,
\]
which can be rewritten as
\[
    -\beta \;\le\; \underbrace{\alpha x - \beta y}_{=:\,r} \;\le\; \alpha.
\]
The invariant $r := \alpha x - \beta y$ therefore takes values in the integer set
\[
    \mathcal{R} \;:=\; \{-\beta,\; -\beta+1,\; \ldots,\; 0,\; \ldots,\; \alpha\},
    \qquad |\mathcal{R}| = \alpha + \beta + 1.
\]

\subsection*{Step 2: Each value of $r$ yields one arithmetic progression}

Recall that the projected quantity used to compute distances is (see equation~(7) of
the companion paper \cite{Kunert2025})
\[
    \tilde{p} \;=\; \beta x + \alpha y.
\]
For a fixed $r \in \mathcal{R}$, consider the linear system
\begin{align}
    \beta x + \alpha y &= \tilde{p}, \label{eq:sys1}\\
    \alpha x - \beta y &= r. \label{eq:sys2}
\end{align}
The coefficient matrix has determinant $-\beta^2 - \alpha^2 = -(\alpha^2+\beta^2)$,
so the unique rational solution is
\[
    x \;=\; \frac{\beta \tilde{p} + \alpha r}{\alpha^2+\beta^2}, \qquad
    y \;=\; \frac{\alpha \tilde{p} - \beta r}{\alpha^2+\beta^2}.
\]
For $x, y \in \mathbb{Z}$ we need $(\alpha^2+\beta^2) \mid (\beta\tilde{p} + \alpha r)$, i.e.
\[
    \tilde{p} \;\equiv\; -\alpha\beta^{-1} r \pmod{\alpha^2+\beta^2},
\]
where $\beta^{-1}$ denotes the modular inverse of $\beta$ modulo $\alpha^2+\beta^2$
(existence is shown in Step~3). Denote this residue class by $c_r$. Then the set
of all $\tilde{p}$ values compatible with invariant $r$ is the arithmetic progression
\[
    P_r \;:=\; \bigl\{\,c_r + k(\alpha^2+\beta^2) \;:\; k \in \mathbb{Z}\,\bigr\}.
\]
Conversely, for every $k \in \mathbb{Z}$ the value $\tilde{p} = c_r + k(\alpha^2+\beta^2)$
yields integer coordinates
\[
    x = \frac{\beta c_r + \alpha r}{\alpha^2+\beta^2} + \beta k, \qquad
    y = \frac{\alpha c_r - \beta r}{\alpha^2+\beta^2} + \alpha k,
\]
both of which are integers by the definition of $c_r$. Hence $P_r \subseteq S$ and
\[
    S \;=\; \bigcup_{r \in \mathcal{R}} P_r.
\]

\subsection*{Step 3: The progressions are pairwise disjoint (Bézout)}

\begin{lemma}
The residues $c_r$ are pairwise distinct modulo $\alpha^2+\beta^2$.
\end{lemma}

\begin{proof}
Suppose $c_{r_1} \equiv c_{r_2} \pmod{\alpha^2+\beta^2}$ for $r_1, r_2 \in \mathcal{R}$.
Then $(\alpha^2+\beta^2) \mid \alpha(r_1 - r_2)$.

By \textbf{Bézout's identity}: $\gcd(\alpha, \beta) = 1$ implies
$\gcd(\alpha, \beta^k) = 1$ for all $k \geq 1$ (any prime dividing both $\alpha$
and $\beta^k$ would divide $\beta$, contradicting $\gcd(\alpha,\beta)=1$).
Therefore
\[
    \gcd\!\bigl(\alpha,\; \alpha^2+\beta^2\bigr) \;=\; \gcd(\alpha,\; \beta^2) \;=\; 1.
\]
Consequently $(\alpha^2+\beta^2) \mid (r_1 - r_2)$.
Since $|r_1 - r_2| \leq \alpha + \beta < \alpha^2 + \beta^2$
(for $(\alpha, \beta) \neq (1,1)$, see Remark~\ref{rem:degenerate}),
we conclude $r_1 = r_2$.
\end{proof}

The same argument gives $\gcd(\beta, \alpha^2+\beta^2) = \gcd(\beta, \alpha^2) = 1$,
so the modular inverse $\beta^{-1}$ required in Step~2 exists.

\subsection*{Step 4: The sorted union has period $\alpha + \beta + 1$}

Sort the $\alpha+\beta+1$ distinct residues in ascending order:
\[
    0 \;\le\; c_{(1)} \;<\; c_{(2)} \;<\; \cdots \;<\; c_{(\alpha+\beta+1)}
    \;<\; \alpha^2+\beta^2.
\]
The sorted union $S$ then reads
\[
    \ldots,\;
    c_{(1)},\; c_{(2)},\; \ldots,\; c_{(\alpha+\beta+1)},\;
    c_{(1)}+D,\; c_{(2)}+D,\; \ldots,\; c_{(\alpha+\beta+1)}+D,\;
    c_{(1)}+2D,\; \ldots
\]
where $D := \alpha^2+\beta^2$. The differences $\delta^{(i)} = \tilde{p}^{(i+1)} - \tilde{p}^{(i)}$
form the pattern
\[
    \underbrace{%
        c_{(2)}-c_{(1)},\;\;
        c_{(3)}-c_{(2)},\;\;
        \ldots,\;\;
        c_{(\alpha+\beta+1)}-c_{(\alpha+\beta)},\;\;
        D - c_{(\alpha+\beta+1)} + c_{(1)}
    }_{\alpha+\beta+1 \text{ terms, summing to } D = \alpha^2+\beta^2}
\]
and this pattern repeats exactly with period $\alpha+\beta+1$.
No smaller period is possible: the $\alpha+\beta+1$ differences are determined by the
$\alpha+\beta+1$ distinct residues, and any sub-period would require a repeated residue,
contradicting Lemma~1.

Therefore $\lambda_{\alpha,\beta} = \alpha + \beta + 1 = \Lambda_{\alpha,\beta}$. \qed

% =====================
\section*{Remarks}
% =====================

\begin{remark}\label{rem:degenerate}
The case $\alpha = \beta = 1$ is degenerate: $\alpha^2+\beta^2 = 2 < \alpha+\beta+1 = 3$,
so three distinct residues modulo~$2$ cannot all differ, and the progressions $P_r$
overlap. In this case the sequence $(\tilde{p})$ contains duplicate values
(e.g.\ both $(x,y) = (0,1)$ and $(1,0)$ project to $\tilde{p} = 1$),
producing zero differences. Nevertheless, one can verify directly that the difference
sequence is $(1, 0, 1, 1, 0, 1, \ldots)$, which has period $3 = \Lambda_{1,1}$.
The theorem therefore holds for $\alpha = \beta = 1$ as well, though the proof
above does not cover this case.
\end{remark}

\begin{remark}
The sum of one full period of differences equals $\alpha^2+\beta^2 = D$,
the common difference of the underlying arithmetic progressions.
This can be used to verify numerical results:
for $(\alpha,\beta) = (2,1)$ the gaps $(1,1,1,2)$ sum to $5 = 4+1$;
for $(\alpha,\beta) = (3,2)$ the gaps $(2,1,2,3,2,3)$ sum to $13 = 9+4$.
\end{remark}

% =========================
\begin{thebibliography}{9}
% =========================

\bibitem{Kunert2025}
Dirk Kunert, 2025,
\emph{Conjectures on the Period Lengths of One-Dimensional Cut-and-Project Sequences},
available at: \url{https://github.com/dkunert/cut-and-project}.

\end{thebibliography}

\end{document}
